% 04_macros.tex
% Componentes Especiales para Instalación y Despliegue (Estándar ISO 15289)

% =============================================================================
% 1. FORMATO PARA COMANDOS, RUTAS Y CADENAS LARGAS
% =============================================================================

% Comando / valor en línea (fondo gris). Cada palabra se encajona por separado
% y se unen con espacios quebrables, de modo que un comando o ruta de menú largo
% pueda partirse entre líneas en lugar de invadir el margen (overfull).
\makeatletter
\newcommand{\comando@chip}[1]{{\setlength{\fboxsep}{2pt}\fontfamily{pcr}\fontsize{9}{10.8}\selectfont\colorbox{grisTecnico}{\texttt{#1}}}}
\def\comando@loop#1 #2\comando@end{%
  \comando@chip{#1}%
  \def\comando@rest{#2}%
  \ifx\comando@rest\@empty\else\space\comando@loop#2\comando@end\fi}
\DeclareRobustCommand{\comando}[1]{\comando@loop#1 \comando@end}
\makeatother

% Rutas, endpoints y rutas de archivo con quiebre automático de línea.
\newcommand{\ruta}[1]{{\ttfamily\fontsize{8.5}{10.2}\selectfont\rutaBreak{#1}}}
\newcommand{\rutaBreak}[1]{\seqsplit{#1}}

% \dir: direcciones/banderas largas (chrome://flags/#..., listas de orígenes)
% que deben romper de forma natural en los separadores sin desbordar el margen.
% Usa \nolinkurl (hyperref + xurl): admite '#' y '_' sin escapar y quiebra en
% '/', ':', '.', '-', etc. Monospace, sin convertirse en hiperenlace.
\newcommand{\dir}[1]{{\fontsize{8.5}{10.2}\selectfont\nolinkurl{#1}}}

% Quiebre lógico tras separadores habituales (/ _ - . @ :) sin perder caracteres.
\makeatletter
\newcommand{\pathbreak}[1]{\texttt{\@pathbreak#1\relax}}
\def\@pathbreak#1{%
  \ifx#1\relax\else
    \ifx#1/\,\discretionary{}{}{}/\discretionary{}{}{}%
    \else\ifx#1_\_\discretionary{}{}{}%
    \else\ifx#1-#1\discretionary{}{}{}%
    \else\ifx#1.#1\discretionary{}{}{}%
    \else\ifx#1:#1\discretionary{}{}{}%
    \else#1\fi\fi\fi\fi\fi
    \expandafter\@pathbreak
  \fi}
\makeatother

% Tipos de columna con CENTRADO VERTICAL del contenido en cada celda.
\newcolumntype{L}[1]{>{\raggedright\arraybackslash}m{#1}}  % izq. + centro vert.
\newcolumntype{C}[1]{>{\centering\arraybackslash}m{#1}}    % centro + centro vert.
\newcolumntype{R}[1]{>{\raggedleft\arraybackslash}m{#1}}   % der. + centro vert.
\newcolumntype{Y}{>{\raggedright\arraybackslash}X}
\newcommand{\celdaCentro}[1]{\multicolumn{1}{c}{#1}}

% =============================================================================
% 2. CAJAS DE AVISO (semántica de despliegue)
% =============================================================================

% Nota: consejos de productividad (fondo azul claro)
\newtcolorbox{AvisoNota}{
    colback=bgNota, colframe=bgNota!80!black,
    boxrule=0.5pt, leftrule=4pt, arc=2pt,
    title={\faInfoCircle\ \textbf{Nota}}, coltitle=black,
    fonttitle=\sffamily, attach title to upper=\quad
}

% Precaución: evitar errores operacionales (borde amarillo)
\newtcolorbox{AvisoPrecaucion}{
    colback=white, colframe=colorAcento,
    boxrule=1.5pt, arc=2pt, fontupper=\bfseries,
    title={\faExclamationTriangle\ \textbf{Precaución}}, coltitle=black,
    fonttitle=\sffamily, attach title to upper=\par
}

% Advertencia: prevenir pérdida de datos (fondo rojo suave)
\newtcolorbox{AvisoAdvertencia}{
    colback=bgAdvertencia, colframe=red!70!black,
    boxrule=0.5pt, leftrule=4pt, arc=2pt,
    title={\faRadiation\ \textbf{Advertencia}}, coltitle=black,
    fonttitle=\sffamily, attach title to upper=\par
}

% Alerta de Seguridad: vulnerabilidades críticas
\newtcolorbox{AlertaSeguridad}{
    colback=bgAdvertencia, colframe=red!80!black,
    boxrule=0.5pt, leftrule=4pt, arc=2pt,
    title={\faShield*\ \textbf{Alerta de Seguridad}}, coltitle=black,
    fonttitle=\sffamily, attach title to upper=\par\vspace{0.2em}
}

% Punto de No Retorno (operación irreversible)
\newtcolorbox{PuntoNoRetorno}{
    colback=bgAdvertencia, colframe=red!80!black,
    boxrule=1.5pt, leftrule=6pt, arc=2pt,
    title={\faRadiation\ \textbf{Punto de No Retorno}}, coltitle=black,
    fonttitle=\sffamily, attach title to upper=\par\vspace{0.2em}
}

% Nota de Configuración (dependencias / efectos secundarios)
\newtcolorbox{NotaConfiguracion}{
    colback=white, colframe=bordeNotaTecnica,
    boxrule=1.2pt, arc=2pt,
    title={\faInfoCircle\ \textbf{Nota de Configuración}}, coltitle=bordeNotaTecnica,
    fonttitle=\sffamily, attach title to upper=\par\vspace{0.2em}
}

% Verificación correcta (resultado esperado tras un paso)
\newtcolorbox{VerificacionOK}{
    colback=bgExito, colframe=colorSecundario!80!black,
    boxrule=0.5pt, leftrule=4pt, arc=2pt,
    title={\faCheckCircle\ \textbf{Resultado Esperado}}, coltitle=black,
    fonttitle=\sffamily, attach title to upper=\par\vspace{0.2em}
}

% =============================================================================
% 3. LISTAS Y CHECKLISTS
% =============================================================================
\newlist{checklist}{itemize}{1}
\setlist[checklist]{label=$\square$, leftmargin=*, noitemsep}

\newlist{ChecklistGoLive}{itemize}{1}
\setlist[ChecklistGoLive]{label=$\square$, leftmargin=*, noitemsep}

% =============================================================================
% 4. ENTORNOS DE TABLA REUTILIZABLES
% =============================================================================

% 4a. Pre-requisitos (Hardware/Software) con casilla de validación
\newenvironment{TablaPrerequisitos}{
    \renewcommand{\arraystretch}{1.5}
    \vspace{0.3cm}\noindent
    \fontsize{9}{10.8}\selectfont
    \tabularx{\dimexpr\textwidth-2pt\relax}{|>{\sffamily\bfseries\color{colorPrimario}}l|Y|c|}
    \hline\rowcolor{headerTabla}
    \sffamily\textbf{Componente} & \sffamily\textbf{Especificación / Versión Mínima} & \sffamily\textbf{Validado} \\
    \hline
}{
    \endtabularx
    \vspace{0.3cm}
}

% 4b. Matriz de Variables de Entorno
\newenvironment{TablaVariables}{
    \renewcommand{\arraystretch}{1.4}
    \vspace{0.3cm}\noindent
    \fontsize{8.5}{10.2}\selectfont\sffamily
    \tabularx{\dimexpr\textwidth-2pt\relax}{|>{\ttfamily\bfseries\color{colorPrimario}\raggedright\arraybackslash}p{4.2cm}|Y|c|}
    \hline\rowcolor{headerTabla}
    \sffamily\textbf{Variable} & \sffamily\textbf{Descripción y origen del valor} & \sffamily\textbf{Oblig.} \\
    \hline
}{
    \endtabularx
    \vspace{0.3cm}
}

% 4c. Tabla genérica de errores (Troubleshooting), multipágina
\newenvironment{TablaErrores}{
    \renewcommand{\arraystretch}{1.35}
    \fontsize{8.5}{10.2}\selectfont\sffamily
    \setlength{\tabcolsep}{4pt}
    \begin{longtable}{|>{\raggedright\arraybackslash}p{3.0cm}|>{\raggedright\arraybackslash}p{3.6cm}|>{\raggedright\arraybackslash}p{8.2cm}|}
    \hline\rowcolor{headerTabla}
    \sffamily\textbf{Síntoma} & \sffamily\textbf{Causa probable} & \sffamily\textbf{Solución} \\
    \hline\endfirsthead
    \hline\rowcolor{headerTabla}
    \sffamily\textbf{Síntoma} & \sffamily\textbf{Causa probable} & \sffamily\textbf{Solución} \\
    \hline\endhead
    \hline\multicolumn{3}{r}{\sffamily\tiny\color{grisISO}\textit{continúa en la página siguiente}}\\
    \endfoot\hline\endlastfoot
}{
    \end{longtable}
}

% 4d. Tabla API (parámetros / endpoints)
\newenvironment{TablaAPI}{
    \renewcommand{\arraystretch}{1.4}
    \vspace{0.3cm}\noindent
    \fontsize{8.5}{10.2}\selectfont
    \tabularx{\dimexpr\textwidth-2pt\relax}{|>{\ttfamily\color{colorPrimario}}l|l|Y|}
    \hline\rowcolor{headerTabla}
    \sffamily\textbf{Endpoint / Parámetro} & \sffamily\textbf{Método} & \sffamily\textbf{Descripción} \\
    \hline
}{
    \endtabularx
    \vspace{0.3cm}
}

% =============================================================================
% 5. BLOQUES DE CÓDIGO Y TERMINAL (Bash, YAML, ENV, JSON, SQL)
% =============================================================================
\lstdefinelanguage{envfile}{
    morekeywords={},
    comment=[l]{\#},
    morestring=[b]",
    sensitive=true
}
\lstdefinestyle{ethosCode}{
    basicstyle=\ttfamily\fontsize{8}{9.6}\selectfont,
    extendedchars=true,
    backgroundcolor=\color{grisTecnico},
    keywordstyle=\color{codeKeyword}\bfseries,
    commentstyle=\color{codeComment}\itshape,
    stringstyle=\color{codeString},
    numbers=left, numberstyle=\tiny\color{grisISO}, numbersep=8pt,
    frame=single, framesep=5pt, rulecolor=\color{grisISO!50},
    breaklines=true, breakatwhitespace=false, showstringspaces=false,
    breakindent=0pt, postbreak=\mbox{\textcolor{grisISO}{$\hookrightarrow$}\space},
    tabsize=2, captionpos=b, aboveskip=0.6em, belowskip=0.6em,
    literate=%
        {á}{{\'a}}1 {é}{{\'e}}1 {í}{{\'\i}}1 {ó}{{\'o}}1 {ú}{{\'u}}1
        {Á}{{\'A}}1 {É}{{\'E}}1 {Í}{{\'I}}1 {Ó}{{\'O}}1 {Ú}{{\'U}}1
        {ñ}{{\~n}}1 {Ñ}{{\~N}}1 {¿}{{?`}}1 {¡}{{!`}}1 {€}{{\euro}}1
        {—}{{-{}-}}2 {–}{{-}}1 {«}{{<{}<}}2 {»}{{>{}>}}2
        {─}{{-}}1 {│}{{|}}1 {•}{{*}}1 {→}{{->}}2 {✓}{{[OK]}}4 {✗}{{[X]}}3
}
\lstset{style=ethosCode}

% Atajo: bloque de terminal etiquetado
\newcommand{\terminal}[1]{\lstinline[basicstyle=\ttfamily\small]|#1|}

% =============================================================================
% 6. ESTILOS TikZ (topologías de red, flujos de despliegue, mockups)
% =============================================================================
\tikzset{
    c4persona/.style={
        rectangle, rounded corners=3pt, draw=c4Persona!80!black, very thick,
        fill=c4Persona, text=white, align=center, font=\sffamily\footnotesize\bfseries,
        minimum width=2.6cm, minimum height=1.2cm, inner sep=4pt},
    c4sistema/.style={
        rectangle, rounded corners=4pt, draw=c4Sistema!80!black, very thick,
        fill=c4Sistema, text=white, align=center, font=\sffamily\footnotesize\bfseries,
        minimum width=3.2cm, minimum height=1.4cm, inner sep=5pt},
    c4contenedor/.style={
        rectangle, rounded corners=4pt, draw=c4Contenedor!70!black, very thick,
        fill=c4Contenedor, text=white, align=center, font=\sffamily\scriptsize\bfseries,
        minimum width=3.0cm, minimum height=1.3cm, inner sep=4pt},
    c4componente/.style={
        rectangle, rounded corners=3pt, draw=c4Componente!70!black, thick,
        fill=c4Componente, text=black, align=center, font=\sffamily\scriptsize,
        minimum width=2.8cm, minimum height=1.0cm, inner sep=3pt},
    c4externo/.style={
        rectangle, rounded corners=4pt, draw=c4Externo!70!black, very thick,
        fill=c4Externo, text=white, align=center, font=\sffamily\footnotesize\bfseries,
        minimum width=2.8cm, minimum height=1.2cm, inner sep=4pt},
    c4datos/.style={
        cylinder, shape border rotate=90, aspect=0.25, draw=c4Datos!70!black, very thick,
        fill=c4Datos, text=white, align=center, font=\sffamily\scriptsize\bfseries,
        minimum width=2.6cm, minimum height=1.4cm, inner sep=4pt},
    c4limite/.style={
        rectangle, rounded corners=6pt, draw=c4BordeLim, dashed, thick, inner sep=10pt},
    c4flecha/.style={-{Stealth[length=2.5mm]}, semithick, draw=grisISO!90!black},
    c4etiqueta/.style={font=\sffamily\scriptsize\itshape, text=grisISO!90!black, align=center,
        fill=white, inner sep=1.5pt},
    flujoInicio/.style={
        rectangle, rounded corners=10pt, draw=flujoInicio!70!black, very thick,
        fill=flujoInicio!85, text=white, align=center, font=\sffamily\scriptsize\bfseries,
        minimum width=2.4cm, minimum height=0.8cm},
    flujoProceso/.style={
        rectangle, draw=c4Sistema!80!black, thick, fill=c4Componente!40,
        align=center, font=\sffamily\scriptsize, minimum width=2.6cm, minimum height=0.9cm,
        inner sep=3pt},
    flujoDecision/.style={
        diamond, aspect=2, draw=flujoDecision!70!black, thick, fill=flujoDecision!30,
        align=center, font=\sffamily\scriptsize, inner sep=1pt, minimum width=2.6cm},
    flujoFin/.style={
        rectangle, rounded corners=10pt, draw=bgFail!70!black, very thick,
        fill=bgFail!85, text=white, align=center, font=\sffamily\scriptsize\bfseries,
        minimum width=2.4cm, minimum height=0.8cm},
    flujoFlecha/.style={-{Stealth[length=2mm]}, semithick}
}
